\section{Terminal junctions and surface centers}\label{sec:junctions}

We now work with a sequence of small diagrams on strong models, and
assume that its initial generalized pair is
generalized klt.  The purpose of this section is to replace the models in
the sequence by crepant generalized terminal models.  Consecutive terminal
models will be joined by a decrease of boundary coefficients followed by a
finite sequence of small elementary steps.  We then show that, after
discarding finitely many diagrams, every coefficient that still changes
belongs to a divisor contracted to a curve or a point.  The terminal-junction
construction follows the method of
\cite[Section~3, Proposition~3.4]{ProjectiveLCFourfolds}; the surface
calculation below supplies the control of coefficient variation needed in
the present analytic setting.

Throughout this section, write the diagrams as
\[
 X_i\xrightarrow{f_i}Z_i\xleftarrow{f_i^+}X_{i+1},
 \qquad D_i=K_{X_i}+B_i+M_{X_i},
\]
and put $a_i(E)=a(E;X_i,B_i+\M)$ for a global place $E$.  We retain all
the hypotheses on the diagrams and on the fixed nef b-divisor from
Theorem~\ref{thm:main}.  In addition, each $X_i$ is strong and
$(X_0,B_0+\M)$ is generalized klt.  The exceptional
loci of $f_i$ and $f_i^+$ are denoted by $L_i$ and $L_i^+$, respectively.
Here the initial boundary and higher nef datum belong to the particular
strong generalized klt sequence under study; a later application may use
the data of a tail obtained after removing a boundary floor.

\begin{proposition}[Crepant terminal junctions]\label{prop:terminal-junctions}
Suppose that the sequence above is infinite.  After passing to a tail,
there are a number $\varepsilon>0$ and a finite set $\mathcal I$ of global
places such that, at every remaining index $i$,
\begin{equation}\label{eq:stable-low-places}
 \begin{gathered}
 a_i(E)\geq\varepsilon\quad\text{for every global place }E,
 \\
 \mathcal I=
 \{E:E\text{ is exceptional over }X_i,\ a_i(E)\leq 1\}.
 \end{gathered}
\end{equation}
There are projective bimeromorphic crepant morphisms
\begin{equation}\label{eq:terminal-junction}
 \begin{gathered}
 h_i:(Y_i,\Theta_i+\M)\longrightarrow (X_i,B_i+\M),
 \\
 K_{Y_i}+\Theta_i+M_{Y_i}=h_i^*D_i,
 \end{gathered}
\end{equation}
with the following properties.
\begin{enumerate}
\item Each $Y_i$ is a strong normal irreducible compact K\"ahler fourfold,
$\Theta_i$ is effective, and
$(Y_i,\Theta_i+\M)$ is generalized terminal.  The exceptional prime
divisors of $h_i$ represent exactly the places in $\mathcal I$.
\item For $E\in\mathcal I$, let $S_{i,E}$ be its prime on $Y_i$.  Define
\begin{equation}\label{eq:junction-coefficient-change}
 \Theta_i'=
 \Theta_i-\sum_{E\in\mathcal I}
       \bigl(a_{i+1}(E)-a_i(E)\bigr)S_{i,E}.
\end{equation}
Starting with $(Y_i,\Theta_i'+\M)$, a finite, possibly empty, sequence
of small elementary steps over $Z_i$ ends at
$(Y_{i+1},\Theta_{i+1}+\M)$.  All its models are strong normal irreducible
compact K\"ahler fourfolds, all its pairs are generalized
terminal, and its boundary and nef data are transported as in an elementary
program.  In particular, every global place is exceptional over one model
of this connecting sequence if and only if it is exceptional over all of
them.
\item There is one finite label set $\mathcal H$ for the primes supplied by
$\mathcal I$ and by the components of $B_0$.  Write $S_h$ for the current
strict transform of a label and $b_h$ for its current coefficient.  Labels
are retained when their coefficient is zero.  On the entire chain of
coefficient changes and small steps,
\[
                    \Theta=\sum_{h\in\mathcal H}b_hS_h,
                    \qquad 0\leq b_h<1,
\]
and every $b_h$ is nonincreasing.  Call a label \emph{varying} if its
coefficient sequence is not eventually constant, and \emph{fixed}
otherwise.  After a further tail is chosen, all fixed coefficients are
constant.  Every varying label comes from $\mathcal I$ and has positive
coefficient at every remaining stage.
\end{enumerate}
\end{proposition}

Figure~\ref{fig:terminal-bridge} records the two kinds of change between
crepant junctions.  The difficulty in Section~\ref{sec:difficulty} will
be compared across both upper moves.
\begin{figure}[ht]
\centering
\begin{tikzcd}[column sep=3.2cm,row sep=large]
 (Y_i,\Theta_i) \arrow[r,Rightarrow,"\Theta_i'\leq\Theta_i"]
   \arrow[d,"h_i"'] &
 (Y_i,\Theta_i') \arrow[r,dashed,"\text{finite small program}"] &
 (Y_{i+1},\Theta_{i+1}) \arrow[d,"h_{i+1}"] \\
 X_i \arrow[r,"f_i"'] & Z_i & X_{i+1}\arrow[l,"f_i^+"]
\end{tikzcd}
\caption{The bridge for one downstairs flip.  The first upper arrow changes
the boundary on the same space; it is not a morphism.  The second denotes a
finite, possibly empty, sequence of small elementary steps over $Z_i$.
The two vertical maps are crepant for the displayed endpoint boundaries,
and the b-divisor $\M$ is fixed throughout.}
\label{fig:terminal-bridge}
\end{figure}

\begin{proof}
By Lemma~\ref{lem:small-comparison}, discrepancies do not decrease, and a
place is exceptional over $X_i$ exactly when it is exceptional over
$X_{i+1}$.  Consequently the sets
\[
 \mathcal I_i=\{E:E\text{ is exceptional over }X_i,\ a_i(E)\leq1\}
\]
form a decreasing sequence.  Proposition~\ref{prop:low-places} makes the
first of them finite and gives a positive lower bound for all initial
discrepancies.  Monotonicity preserves that bound.  A decreasing sequence
of subsets of a finite set stabilizes, which gives
\eqref{eq:stable-low-places}.  We reindex this tail.

If $\mathcal I$ is empty, take the identity as the first junction.
Otherwise apply Proposition~\ref{prop:extraction} at the first index to
extract exactly $\mathcal I$.  Its hypotheses hold: the model is strong,
the pair is generalized klt with rational effective boundary, and the
fixed nef data are carried on a higher model projective over $X_i$.
The resulting morphism is projective bimeromorphic and has a strong normal
irreducible compact K\"ahler source.  Crepancy gives the coefficient
$1-a_i(E)\in[0,1-\varepsilon]$ on an extracted prime; all other boundary
components are strict transforms of $B_i$.  Thus the crepant boundary is
effective.  Notice that a place with $a_i(E)=1$ is extracted with coefficient
zero.

Any place exceptional over this source is exceptional over $X_i$: a prime
on $X_i$ has a strict transform on every proper birational model.  Such a
place is not one of the extracted primes, so its discrepancy is strictly
greater than $1$ by \eqref{eq:stable-low-places}.  Crepancy and the positive
lower bound prove generalized terminality.  This constructs the first
junction.

Suppose the junction over $X_i$ has been constructed.  The coefficient of
$S_{i,E}$ changes in \eqref{eq:junction-coefficient-change} from
$1-a_i(E)$ to $1-a_{i+1}(E)$.  Both are nonnegative because $E$ belongs to
the stabilized set at both indices, and the latter is no larger than the
former.  The divisor subtracted from $\Theta_i$ is effective and
$\Q$-Cartier on the strong model $Y_i$.  Its pullback to a higher model is
effective; hence subtracting it only increases discrepancies.  In
particular, $\Theta_i'$ remains effective and
$(Y_i,\Theta_i'+\M)$ remains generalized terminal.  Set
\[
                    D_i'=K_{Y_i}+\Theta_i'+M_{Y_i}.
\]

We first compare this adjoint with the positive model downstairs.  By
Lemma~\ref{lem:common-projective-models}, choose a smooth common
model $W$, carrying $\M$, with projective maps
$q:W\to Y_i$ and $r:W\to X_{i+1}$, all over $Z_i$, and put
\begin{equation}\label{eq:junction-positive-comparison}
                         F=q^*D_i'-r^*D_{i+1}.
\end{equation}
The coefficient change makes $q_*F=0$.  Indeed, at a prime $S_{i,E}$ with
$E\in\mathcal I$, the coefficient of the crepant boundary computed from
$X_{i+1}$ is $1-a_{i+1}(E)$, the new coefficient in $\Theta_i'$.
Every other prime on $Y_i$ corresponds, by smallness downstairs, to a
prime on $X_{i+1}$ and has its transported boundary coefficient.  The
canonical representatives and the trace of the same b-divisor agree in
this calculation.  On a curve contracted by $q$, the divisor $-F$ has
nonnegative degree: $q^*D_i'$ has degree zero and $r^*D_{i+1}$ is pulled
back from the $f_i^+$-ample divisor $D_{i+1}$.  The negativity
lemma therefore gives $F\geq0$.  Moreover every prime on $X_{i+1}$ is
represented on $Y_i$, where the coefficient of $F$ is zero.  Thus $F$ is
exceptional over $X_{i+1}$.

Use the finite generalized klt part of
Proposition~\ref{prop:relative-program} for $D_i'$ over $Z_i$.  The map
$Y_i\to Z_i$ is projective bimeromorphic to a normal compact K\"ahler
base, the source is strong and generalized klt, and a common higher model
projective over $Y_i$ carries the fixed data with the
required relative nefness.  The proposition supplies a finite elementary
program ending at a model $\overline Y$ whose adjoint $\overline D$ is nef
over $Z_i$.  On a further common model, write again $q$ for the map to
$Y_i$ and $\bar q$ for the map to $\overline Y$.
Lemma~\ref{lem:program-comparison} gives
\[
              q^*D_i'=\bar q^*\overline D+G,
              \qquad G\geq0\text{ exceptional over }\overline Y.
\]
Together with \eqref{eq:junction-positive-comparison}, these are the two
effective comparisons required by Lemma~\ref{lem:nef-ample}.  That lemma
gives a projective morphism $h_{i+1}:\overline Y\to X_{i+1}$ and the
actual crepant identity $\overline D=h_{i+1}^*D_{i+1}$.

This program cannot have lost a marked prime $S_{i,E}$.  Its discrepancy
immediately after the coefficient change is $a_{i+1}(E)$; by the crepant
identity its discrepancy at the end is again $a_{i+1}(E)$.  If a divisorial
step had contracted it, the strict comparison for that step would have
increased this discrepancy, and subsequent comparisons could not decrease
it.  Nor can the program have lost any other prime on $Y_i$.  Such a prime
corresponds to a prime on $X_{i+1}$, and that prime must have a strict
transform on $\overline Y$.  An elementary program extracts no prime, so a
prime once contracted cannot reappear.  We have excluded every divisorial
step.  All connecting steps are therefore small, and their discrepancy
comparisons preserve generalized terminality.  They also preserve the
exceptionality status of every place.  Smallness downstairs keeps each
marked place exceptional over $X_{i+1}$, whereas every other surviving
prime corresponds to a prime on $X_{i+1}$.  Thus the surviving marked
primes are exactly the exceptional primes of $h_{i+1}$, and the end is the
required junction $Y_{i+1}$.  This proves the construction by induction.  The
relative program is allowed to have length zero.

There are finitely many marked primes and finitely many components of
$B_0$.  Smallness transports all of them through every connector.  The
downstairs coefficients are constant, and the marked coefficients are
$1-a_i(E)$ at junctions and are changed only by
\eqref{eq:junction-coefficient-change}.  They are therefore nonincreasing
and lie in $[0,1)$.  For the finitely many coefficient sequences that are
eventually constant, discard a prefix after all have become constant.
Every other sequence stays positive: a nonnegative nonincreasing sequence
that reaches zero is thereafter zero.  Only marked coefficients can be of
this latter type.  This proves the last assertion.
\end{proof}

All the terminal models just constructed are also ordinarily terminal by
Lemma~\ref{lem:ordinary-singularities}.  In particular they are smooth in
codimension two.  To control marked coefficients above a surface, we need
no bound on the number of exceptional curves in a transverse slice.
Grouping the curves by their global marked labels will give a linear
system with at most $|\mathcal I|$ unknowns and an integer matrix whose
columns are linearly independent.  We will bound its entries and clear one denominator for
its right-hand side; a nonsingular minor then gives a uniform lattice for
the marked coefficients.

\begin{proposition}[Finiteness above marked surfaces]\label{prop:surface-finiteness}
In the setting of Proposition~\ref{prop:terminal-junctions}, there is a
finite set
\[
                    \mathcal A_{\mathrm{surf}}
                    \subset [\varepsilon,1]\cap\Q
\]
such that $a_i(E)\in\mathcal A_{\mathrm{surf}}$ whenever
$E\in\mathcal I$ and $\cent_{X_i}(E)$ is a surface.  After passing to a
further tail, the following conclusions hold.
\begin{enumerate}
\item No $E\in\mathcal I$ has a surface center contained in $L_i$ on
$X_i$, or a surface center contained in $L_i^+$ on $X_{i+1}$.
\item Every varying label $S_h\subset Y_i$ satisfies
$\dim h_i(S_h)\leq1$ at every junction.
\item For an irreducible compact analytic surface contained in either $L_i$
or $L_i^+$, the corresponding
junction morphism is an isomorphism near its general point, and the
downstairs model is smooth there.  Blowing up the coherent ideal of that
surface defines a unique global generic blowup place.  It is exceptional
over the corresponding downstairs model and has discrepancy strictly
greater than $1$; in particular it does not belong to $\mathcal I$.
\end{enumerate}
\end{proposition}

\begin{proof}
We first prove the assertion about $\mathcal A_{\mathrm{surf}}$.  If
$\mathcal I$ is empty, choose $\mathcal A_{\mathrm{surf}}=\varnothing$.
For the remaining case, fix a junction $\rho:Y\to X$ and a surface
$V\subset X$ which is the center of at least one marked prime.  We work
near general points of $V$.

\medskip\noindent\emph{A transverse surface and its curves.}
The exceptional image of a proper birational morphism to normal $X$ has
codimension at least two.  We may avoid every other component of that
image and every center properly contained in $V$.  The full inverse image
of $V$ is a proper analytic subset of the irreducible fourfold $Y$, hence
has dimension at most three.  A component supplying a curve over a general
point of $V$ must therefore have dimension three and be a divisor
dominating $V$.  All such divisors are marked.  Components of dimension
at most two cannot supply a curve in a general fiber.  These observations
also show that the general fibers over $V$ have dimension at most one.

The singular locus of terminal $Y$ has dimension at most one, so it misses
the fiber over a general point of $V$.  The singular loci of the marked
divisors, their pairwise intersections, and their intersections with
strict transforms of components of $B_0$ have dimension at most two.
Choose also a projective carrier $\sigma:W\to Y$ for the fixed nef data,
with $W$ smooth.  Its nonisomorphism image has dimension at most two.  None of
these subsets can contain an entire curve over a general point of $V$:
a component that dominates the surface has zero-dimensional general
fiber, and components with smaller image can be avoided.

Here is a precise choice of the slice used below.  Near a general smooth
point of $V$, lift two local coordinates on $V$ to holomorphic functions
on $X$.  Take a general common level set of their pullbacks on $Y$.
Finite analytic stratifications near the proper fiber can be chosen to
respect the divisors and the subsets in the preceding paragraph.  On
every stratum dominating $V$, generic smoothness makes the two pulled-back
functions submersive over general values; the images of all other strata
can be avoided.  It follows, after shrinking about the fiber, that the
level set is a smooth surface $R$.  The restriction of $K_Y$ to $R$ is
$K_R$, because the two slice equations trivialize its normal bundle.

Let $C_1,\ldots,C_q$ be the reduced irreducible curves of the fiber in
$R$.  Each is supplied by exactly one marked divisor, with multiplicity
one in that divisor's slice: at a general point of the curve the divisor
is smooth, the slice is transverse on it, and none of the other listed
subsets contains that curve.  Write $\lambda(j)$ for its label.  The same
label may supply several curves, but each marked divisor with center $V$
supplies at least one.  No $C_j$ is contained in a strict downstairs
boundary component or in the nonisomorphism image of $\sigma$.

The restricted map on the slice is proper near this fiber, contracts the
$C_j$ to a point, and is an isomorphism off that point locally.  No
surface component is mapped to the point, by the fiber dimension bound.
Pass, if necessary, to the normal Stein factor of the local target, or
equivalently to the normalization of the relevant local image.  We then
have a proper modification from the smooth surface whose reduced
exceptional set is the union of these curves.  Grauert's
negative-definiteness criterion
\cite[Section~4, no.~8(e), pp.~366--367]{Grauert1962} therefore applies:
the matrix
\[
                       A=(C_j\cdot C_k)_{1\leq j,k\leq q}
\]
is negative definite.  Put $u_j=-C_j^2\in\Z_{>0}$ and
$m_{jk}=C_j\cdot C_k\in\Z_{\geq0}$ for $j\ne k$.

\medskip\noindent\emph{Uniform bounds for the intersection equations.}
Let $B_Y^{\mathrm{str}}$ denote the strict boundary part coming from
$B_i$.  It has nonnegative intersection with every $C_j$, since it is
effective and contains no such curve.  The same nonnegativity holds for
$M_Y$.  Indeed, the strict transform $C_j'$ on $W$ meets the isomorphism
locus of $\sigma$.  For the effective nef defect of
Lemma~\ref{lem:nef-defect},
\[
                         J=\sigma^*M_Y-M_W\geq0,
\]
the curve $C_j'$ is not contained in $\Supp J$.  The projection formula,
nefness on the carrier, and effectivity then give
\[
                 M_Y\cdot C_j=M_W\cdot C_j'+J\cdot C_j'\geq0.
\]
Intersect the crepant equation
$K_Y+\Theta+M_Y=\rho^*D_X$ with $C_j$.  On $R$ the marked part near the
fiber is $\sum_k b_{\lambda(k)}C_k$.  Adjunction for the reduced
irreducible compact curve $C_j$ gives
$K_R\cdot C_j=2p_a(C_j)-2+u_j$.  We obtain
\begin{align*}
 0={}&2p_a(C_j)-2+(1-b_{\lambda(j)})u_j
       +\sum_{k\ne j}b_{\lambda(k)}m_{jk}
       +(B_Y^{\mathrm{str}}+M_Y)\cdot C_j\\
 \geq{}&2p_a(C_j)-2+(1-b_{\lambda(j)})u_j.
\end{align*}
Since $p_a(C_j)\geq0$ and
$1-b_{\lambda(j)}\geq\varepsilon$, it follows that
\begin{equation}\label{eq:surface-self-intersection-bound}
 (1-b_{\lambda(j)})u_j\leq2-2p_a(C_j)\leq2,
 \qquad u_j\leq U:=\left\lceil\frac2\varepsilon\right\rceil.
\end{equation}

Negative definiteness now bounds each whole row of $A$, even when $q$ is
large.  Fix $j$ and test its quadratic form on the vector whose $j$th
entry is $1$ and whose $k$th entry is $m_{jk}/u_k$ for $k\ne j$.  Its
value is
\[
 -u_j+\sum_{k\ne j}\frac{m_{jk}^2}{u_k}
 +2\sum_{\substack{k<\ell\\k,\ell\ne j}}
   \frac{m_{k\ell}m_{jk}m_{j\ell}}{u_ku_\ell}<0.
\]
The last sum is nonnegative.  Therefore
\begin{equation}\label{eq:surface-row-bound}
 \sum_{k\ne j}\frac{m_{jk}^2}{u_k}<u_j,
 \qquad
 \sum_{k\ne j}m_{jk}\leq\sum_{k\ne j}m_{jk}^2<U^2.
\end{equation}
The second inequality uses $u_k\leq U$ and the integrality and
nonnegativity of every $m_{jk}$.  The absolute sum in each row of $A$ is
thus at most $C_{\mathrm{row}}:=U+U^2$.

\medskip\noindent\emph{Grouping curves by global labels.}
Let $\mathcal H_V$ be the marked labels whose center is $V$, and put
$t=|\mathcal H_V|\leq|\mathcal I|$.  Define the $q$-by-$t$ zero-one
matrix $P$ by
\[
             P_{k,h}=\begin{cases}1,&\lambda(k)=h,\\0,&\lambda(k)\ne h.
                         \end{cases}
\]
Its columns have disjoint nonempty supports, so $P$ has column rank $t$.
Since $A$ is invertible, $AP$ has the same column rank.  Each of its
integer entries is a sum from one row of $A$, and hence has absolute
value at most $C_{\mathrm{row}}$.  The intersection equations take the
form
\begin{equation}\label{eq:grouped-surface-equations}
 \begin{gathered}
                  (AP)(b_h)_{h\in\mathcal H_V}=-(v_j)_{1\leq j\leq q},
 \\
 v_j=K_R\cdot C_j+(B_Y^{\mathrm{str}}+M_Y)\cdot C_j.
 \end{gathered}
\end{equation}
No marked coefficient occurs in the right-hand side.

There is a fixed integer $m_{\mathrm{data}}>0$ such that every
$m_{\mathrm{data}}v_j$ is integral, for every junction and surface in
this calculation.  Choose it to clear the initial boundary coefficients
of this strong generalized klt sequence and to make its chosen higher nef
divisor integral Cartier.  Canonical
degrees on the smooth $R$ are integral.  Each strict boundary prime is
Cartier in the smooth neighborhood of the fiber, so its degree has the
chosen boundary denominator.  Finally, the trace of the integral Cartier
divisor $m_{\mathrm{data}}\M$ on any model is an integral Weil divisor:
compute it by pulling back to a common higher model and pushing forward.
On the smooth neighborhood of the present fiber that trace is Cartier,
so its degree is integral as well.  No assertion about its Cartier index
elsewhere is needed.

Choose $t$ rows of $AP$ giving a nonsingular square matrix.  Its
determinant is a nonzero integer with an absolute value bounded in terms
of $t$ and $C_{\mathrm{row}}$, for example by
$t^{t/2}C_{\mathrm{row}}^t$ by Hadamard's inequality.  Since
$1\leq t\leq|\mathcal I|$, choose one integer $H_{\mathrm{det}}$ bounding
all these absolute values.  Cramer's rule applied to
\eqref{eq:grouped-surface-equations} then gives
\[
 \begin{gathered}
 q_{\mathrm{surf}}=
 m_{\mathrm{data}}\operatorname{lcm}(1,2,\ldots,H_{\mathrm{det}}),
 \\
 b_h\in q_{\mathrm{surf}}^{-1}\Z\qquad(h\in\mathcal H_V).
 \end{gathered}
\]
The denominator is independent of the junction, the surface, and the
number of geometric curves in the slice.  In particular the marked
discrepancies above surfaces belong to the finite set
\begin{equation}\label{eq:marked-surface-discrepancies}
          \mathcal A_{\mathrm{surf}}
          =[\varepsilon,1]\cap q_{\mathrm{surf}}^{-1}\Z.
\end{equation}
This finite set is used only to exclude recurring marked surface centers
contained in the exceptional loci.  The weight denominator in
Section~\ref{sec:difficulty} will be chosen later from the actual fixed
coefficients and the initial rational data; it need not clear every element
of $\mathcal A_{\mathrm{surf}}$.

\medskip\noindent\emph{Removing marked surface centers from the exceptional loci.}
Fix $E\in\mathcal I$.  Suppose its center is a surface in $L_i$ at
infinitely many indices $i_1<i_2<\cdots$.  The strict comparison at
$i_k$ and monotonicity between indices give
\[
                  a_{i_{k+1}}(E)\geq a_{i_k+1}(E)>a_{i_k}(E).
\]
Both $a_{i_k}(E)$ and $a_{i_{k+1}}(E)$ belong to
$\mathcal A_{\mathrm{surf}}$, a contradiction.  If instead its center is
a surface in $L_i^+$ at
infinitely many indices, the corresponding positive-side discrepancies
satisfy
\[
             a_{i_{k+1}+1}(E)>a_{i_{k+1}}(E)\geq a_{i_k+1}(E),
\]
again contradicting membership of $a_{i_k+1}(E)$ and
$a_{i_{k+1}+1}(E)$ in the same finite set.  There are finitely many
$E\in\mathcal I$, so one tail excludes all
such occurrences on both sides.  This proves assertion (1).

On that tail, suppose a marked prime has a surface center $V$ at a
junction over $X_i$.  Assertion (1) says $V$ is not contained in $L_i$.
Its general point is therefore in the common isomorphism open of the
next diagram.  The next center is its strict transform, still a surface,
and Lemma~\ref{lem:small-comparison} gives equality of its two
discrepancies.  Repeating the argument at every subsequent index shows
that its coefficient is constant from this index onward.  A varying label
cannot have such a surface center.  Its center cannot be a divisor
downstairs because it represents an exceptional place.  This proves
assertion (2).

Finally let $V$ be an irreducible compact analytic surface in one of the exceptional loci of a
remaining diagram, and use the junction on that side.  If the junction
were not an isomorphism at the general point of $V$, its general fiber
there would be positive dimensional: a proper birational morphism to a
normal space is an isomorphism wherever it is finite.  A component
supplying those curves over $V$ would be a divisor with center $V$, by
the dimension argument at the start of the proof.  This is excluded by
assertion (1).  Thus the junction is an isomorphism near a general point
of $V$.  The source is ordinarily terminal and smooth in codimension two,
so the downstairs model is smooth at a general point of this surface.

The coherent ideal of the global analytic surface $V$ has a projective
blowup.  Over the dense open where the ambient model and $V$ are smooth,
this is the ordinary codimension-two blowup and has a unique prime
divisor dominating $V$.  Its closure on the normalized blowup defines
a global place with center $V$, hence an exceptional place on the
downstairs model.  If its discrepancy were at most $1$, it would belong
to the stabilized set $\mathcal I$, contradicting assertion (1).
This proves assertion (3).
\end{proof}

At each remaining junction, the reduced union of the varying labels is
contracted by $h_i$ to a set of dimension at most one.  This is the input
for Theorem~\ref{thm:branch}.  Assertion~(3) has a separate later use:
over a surface in the positive exceptional locus of a downstairs diagram,
it supplies the global generic blowup place and the smooth downstairs
neighborhood used by the detector in Section~\ref{sec:completion}.
